\part{General Equilibrium and the State}

\chapter{General Equilibrium and Welfare Economics}
\label{chap:general-equilibrium}

\section{Learning objectives}

After completing this chapter, you should be able to:
\begin{itemize}
  \item distinguish partial from general equilibrium;
  \item explain the Edgeworth box and Pareto efficiency;
  \item state the two welfare theorems with their assumptions; and
  \item distinguish efficiency from equity.
\end{itemize}

\section{From partial to general equilibrium}

Partial-equilibrium analysis studies one market while treating conditions in other markets as fixed. \keyterm{General equilibrium} recognises that markets are connected. A change in the price of electricity affects household budgets, firms' costs, demand for appliances, wages in energy-intensive sectors, and public revenue. General-equilibrium analysis asks whether all these choices are mutually consistent at the same time.

\section{Exchange and the Edgeworth box}

Consider two people, Ana and Bruno, and two goods, $x$ and $y$, with fixed total amounts. An \keyterm{Edgeworth box} represents every feasible allocation. Its width is the total amount of $x$ and its height is the total amount of $y$. Each point gives Ana's bundle from one origin and Bruno's remaining bundle from the opposite origin.

An allocation is \keyterm{Pareto efficient} if no feasible change can make at least one person better off without making anyone worse off. At an interior efficient allocation, both people have the same marginal rate of substitution:
\[
MRS^A_{x,y}=MRS^B_{x,y}.
\]
Otherwise, they could exchange goods in a way that benefits both.

\section{Competitive equilibrium}

A competitive equilibrium is a set of prices and allocations such that every consumer maximises utility subject to their budget and every firm maximises profit, while all markets clear. Prices coordinate dispersed decisions: they communicate scarcity and create incentives without requiring a central planner to know each person's preferences or each firm's technology.

\section{The welfare theorems}

The \keyterm{First Welfare Theorem} states that, under demanding assumptions, every competitive equilibrium is Pareto efficient. Key assumptions include price-taking behaviour, complete markets, well-defined property rights, and the absence of externalities or significant information failures.

The \keyterm{Second Welfare Theorem} states, again under strong assumptions, that any Pareto-efficient allocation can be decentralised as a competitive equilibrium after an appropriate redistribution of initial resources. It separates two questions: society can choose a distribution through transfers, then allow markets to allocate resources efficiently.

\section{Efficiency is not equity}

Pareto efficiency does not say that an allocation is fair. A society in which one person owns nearly everything can be Pareto efficient if any transfer makes that person worse off. The welfare theorems therefore do not settle questions about equality, rights, needs, or political legitimacy. They show conditions under which markets can coordinate efficiently; they do not provide a complete theory of justice.

\section*{Chapter Review}

\subsection*{Key concepts}
\begin{itemize}
  \item \textbf{General equilibrium}: simultaneous equilibrium in interconnected markets.
  \item \textbf{Pareto efficiency}: no feasible Pareto improvement remains.
  \item \textbf{Competitive equilibrium}: optimisation by agents plus market clearing.
  \item \textbf{First Welfare Theorem}: competitive equilibrium is efficient under specified assumptions.
  \item \textbf{Second Welfare Theorem}: redistribution can separate equity from efficient allocation under specified assumptions.
\end{itemize}

\section*{Exercises}

\subsection*{Conceptual questions}
\begin{enumerate}[label=\arabic*.]
  \item Distinguish partial and general equilibrium using an example.
  \item Define a Pareto improvement and Pareto efficiency.
  \item Why must MRS values be equal at an interior efficient exchange allocation?
  \item List three assumptions needed for the First Welfare Theorem.
\end{enumerate}

\subsection*{Graphical exercises}
\begin{enumerate}[resume,label=\arabic*.]
  \item Draw an Edgeworth box and identify a feasible allocation, a Pareto improvement, and a contract curve.
  \item Use indifference curves to show why an allocation with unequal MRS values is not efficient.
\end{enumerate}

\subsection*{Analytical exercises}
\begin{enumerate}[resume,label=\arabic*.]
  \item Ana has $u_A(x_A,y_A)=x_Ay_A$ and Bruno has $u_B(x_B,y_B)=x_By_B$. Write the condition for an interior Pareto-efficient allocation.
  \item Explain why an allocation can be Pareto efficient while highly unequal.
\end{enumerate}

\subsection*{Exam-style multiple-choice questions}
\begin{enumerate}[resume,label=\arabic*.]
  \item An allocation is Pareto efficient when: (a) everyone has equal income; (b) no Pareto improvement is feasible; (c) the government does not intervene; or (d) firms earn zero profit.
  \item The First Welfare Theorem requires, among other conditions: (a) externalities; (b) price-taking behaviour; (c) monopoly; or (d) incomplete markets.
  \item The Second Welfare Theorem highlights that: (a) equity and efficiency are identical; (b) redistribution may be considered separately from market allocation; (c) all redistributions are costless in reality; or (d) prices are unnecessary.
\end{enumerate}

\subsection*{PPE discussion question}

\noindent Do the welfare theorems support a political philosophy of laissez-faire? Discuss which of their assumptions and which concerns about justice must be examined before drawing that conclusion.
